\documentclass[tikz,border=6pt]{standalone}
\usepackage[UTF8]{ctex}
\usetikzlibrary{arrows.meta,calc,positioning}

\definecolor{tsgreen}{HTML}{208443}
\definecolor{tsgreenbg}{HTML}{EEF8F2}
\definecolor{tsorange}{HTML}{E7A82A}
\definecolor{tsorangebg}{HTML}{FFF7E8}
\definecolor{tsred}{HTML}{BE3431}
\definecolor{tsredbg}{HTML}{FFF1F2}
\definecolor{tsviolet}{HTML}{7A59C7}
\definecolor{tsvioletbg}{HTML}{F6F1FF}
\definecolor{tsblue}{HTML}{2F6FDA}
\definecolor{tsbluebg}{HTML}{F2F6FF}
\definecolor{tsborder}{HTML}{D9E2EE}
\definecolor{tstext}{HTML}{1D2738}
\definecolor{tsmuted}{HTML}{6C7A90}

\tikzset{
  >={Latex[length=2.8mm,width=2.3mm]},
  lane/.style={draw=tsborder, rounded corners=8pt, line width=0.8pt, fill=white},
  head/.style={
    rounded corners=6pt, line width=0.75pt, fill=white,
    minimum height=0.55cm, inner sep=4pt, align=center
  },
  stepcircle/.style={
    circle, draw=tsblue!35, line width=0.8pt, fill=tsbluebg,
    minimum size=0.58cm, inner sep=0pt, font=\bfseries\fontsize{9.2}{10.5}\selectfont, text=tsblue
  },
  card/.style={
    draw=tsborder, rounded corners=7pt, line width=0.75pt,
    fill=white, inner sep=5pt, align=left
  },
  wcard/.style={card, fill=tsgreenbg, draw=tsgreen!35!black},
  hcard/.style={card, fill=tsorangebg, draw=tsorange!70},
  gcard/.style={card, fill=tsredbg, draw=tsred!30},
  scard/.style={card, fill=tsvioletbg, draw=tsviolet!35},
  rcard/.style={card, fill=tsbluebg, draw=tsblue!35},
  footG/.style={draw=tsgreen!40!black, fill=tsgreenbg, line width=0.75pt, rounded corners=6pt, inner sep=5pt, font=\fontsize{9.2}{11.5}\selectfont, text=tsgreen!75!black, align=center},
  footR/.style={draw=tsred!30, fill=tsredbg, line width=0.75pt, rounded corners=6pt, inner sep=5pt, font=\fontsize{9.2}{11.5}\selectfont, text=tsred!80!black, align=center},
  footO/.style={draw=tsorange!70, fill=tsorangebg, line width=0.75pt, rounded corners=6pt, inner sep=5pt, font=\fontsize{9.2}{11.5}\selectfont, text=tsorange!85!black, align=center},
  footB/.style={draw=tsblue!35, fill=tsbluebg, line width=0.75pt, rounded corners=6pt, inner sep=5pt, font=\fontsize{9.2}{11.5}\selectfont, text=tsblue, align=center},
  gflow/.style={draw=tsgreen, line width=1.0pt, -{Latex[length=2.9mm,width=2.7mm]}},
  rflow/.style={draw=tsred, line width=1.0pt, dashed, -{Latex[length=2.9mm,width=2.7mm]}},
  bflow/.style={draw=tsblue, line width=1.0pt, -{Latex[length=2.9mm,width=2.7mm]}},
}

\newcommand{\stitle}[1]{{\bfseries\fontsize{11.0}{12.5}\selectfont #1}}
\newcommand{\sbody}[1]{{\fontsize{9.5}{12.8}\selectfont #1}}

\begin{document}
\begin{tikzpicture}[x=1cm,y=1cm]

\node[anchor=west, font=\bfseries\fontsize{18}{20}\selectfont, text=tstext] at (0,11.95)
  {图2-2\quad 端到端软件流转时序图};
\node[anchor=west, font=\fontsize{10.6}{13}\selectfont, text=tsmuted] at (0,11.42)
  {流程与第三章 3.2 小节保持一致：八步闭环、两类边界、一个预 SPU 权威门。};

% step circles
\foreach \i/\y in {1/10.08,2/9.00,3/7.92,4/6.84,5/5.76,6/4.68,7/3.60,8/2.52}{
  \node[stepcircle] at (-0.45,\y) {\i};
}

% lane frames
\draw[lane] (-0.02,1.40) rectangle (2.65,10.92);
\draw[lane] (2.95,1.40) rectangle (5.62,10.92);
\draw[lane] (5.92,1.40) rectangle (8.59,10.92);
\draw[lane] (8.89,1.40) rectangle (11.56,10.92);
\draw[lane] (11.86,1.40) rectangle (14.53,10.92);
\draw[lane] (14.83,1.40) rectangle (17.50,10.92);

% heads
\node[head, draw=tsblue!35, fill=tsbluebg, text=tstext, font=\bfseries\fontsize{12}{13}\selectfont, text width=2.05cm] at (1.315,10.48) {用户 / 主线程};
\node[head, draw=tsgreen!35!black, fill=tsgreenbg, text=tstext, font=\bfseries\fontsize{12}{13}\selectfont, text width=2.05cm] at (4.285,10.48) {浏览器 Worker};
\node[head, draw=tsorange!70, fill=tsorangebg, text=tstext, font=\bfseries\fontsize{12}{13}\selectfont, text width=2.05cm] at (7.255,10.48) {HTTP / multipart};
\node[head, draw=tsred!30, fill=tsredbg, text=tstext, font=\bfseries\fontsize{12}{13}\selectfont, text width=2.05cm] at (10.225,10.48) {服务端快检};
\node[head, draw=tsviolet!35, fill=tsvioletbg, text=tstext, font=\bfseries\fontsize{12}{13}\selectfont, text width=2.05cm] at (13.195,10.48) {SPU / 安全执行};
\node[head, draw=tsblue!35, fill=tsbluebg, text=tstext, font=\bfseries\fontsize{12}{13}\selectfont, text width=2.05cm] at (16.165,10.48) {审计与回显};

% guide lines
\foreach \x in {1.315,4.285,7.255,10.225,13.195,16.165}{
  \draw[tsborder, line width=0.55pt] (\x,1.78) -- (\x,9.82);
}

% user lane
\node[card, text width=1.75cm, minimum height=1.25cm, anchor=north] (choose) at (1.315,10.05) {%
  \stitle{加载本地影像}\\[0.08cm]
  \sbody{主线程保留文件句柄与页面状态。}
};
\node[card, text width=1.75cm, minimum height=1.00cm, anchor=north] (show) at (1.315,2.30) {%
  \stitle{页面展示}\\[0.08cm]
  \sbody{主线程仅展示分类、质量与审计结果。}
};

% worker lane
\node[wcard, text width=1.78cm, minimum height=0.92cm, anchor=north] (recv) at (4.285,10.05) {%
  \stitle{接收句柄}\\[0.06cm]
  \sbody{确认单活 Worker，准备本地计算。}
};
\node[wcard, text width=1.78cm, minimum height=1.05cm, anchor=north] (sniff) at (4.285,8.55) {%
  \stitle{头部尺寸嗅探}\\[0.06cm]
  \sbody{解析 PNG/JPEG/WebP 头部，优先阻断像素炸弹。}
};
\node[wcard, text width=1.78cm, minimum height=0.95cm, anchor=north] (dqa) at (4.285,7.18) {%
  \stitle{解码与预检}\\[0.06cm]
  \sbody{中心裁剪、归一化与 DQA 指标提取。}
};
\node[wcard, text width=1.78cm, minimum height=1.03cm, anchor=north] (share) at (4.285,5.85) {%
  \stitle{share 与哈希}\\[0.06cm]
  \sbody{生成两份 share、一次性标识与审计链。}
};

% http lane
\node[hcard, text width=1.66cm, minimum height=0.92cm, anchor=north] (multipart) at (7.255,5.85) {%
  \stitle{multipart 提交}\\[0.06cm]
  \sbody{跨域仅发送 share 与结构化摘要。}
};

% guard lane
\node[gcard, text width=1.78cm, minimum height=0.92cm, anchor=north] (proto) at (10.225,4.95) {%
  \stitle{协议预检}\\[0.06cm]
  \sbody{长度门、boundary、字段集与 JSON 字节门。}
};
\node[gcard, text width=1.78cm, minimum height=0.98cm, anchor=north] (tensor) at (10.225,3.70) {%
  \stitle{张量快检}\\[0.06cm]
  \sbody{share 哈希、NaN/Inf、次正规数与 DQA 复核。}
};

% spu lane
\node[scard, text width=1.78cm, minimum height=0.86cm, anchor=north] (enter) at (13.195,3.70) {%
  \stitle{进入 SPU}\\[0.06cm]
  \sbody{只有通过权威快检后才触发 2PC 前向。}
};
\node[scard, text width=1.78cm, minimum height=0.92cm, anchor=north] (infer) at (13.195,2.45) {%
  \stitle{密态前向}\\[0.06cm]
  \sbody{动态安全剪枝与分类决策在安全图内部完成。}
};

% response lane
\node[rcard, text width=1.78cm, minimum height=0.90cm, anchor=north] (audit) at (16.165,2.45) {%
  \stitle{审计落盘}\\[0.06cm]
  \sbody{记录载荷指纹、质量裁决与控制面指标。}
};
\node[rcard, text width=1.78cm, minimum height=0.90cm, anchor=north] (reply) at (16.165,1.35) {%
  \stitle{结果回显}\\[0.06cm]
  \sbody{仅回传分类结论、质量裁决与审计摘要。}
};

% flows
\draw[gflow] (choose.east) -- (recv.west);
\draw[gflow] ($(recv.south)+(0,0.01)$) -- ($(sniff.north)+(0,-0.01)$);
\draw[gflow] ($(sniff.south)+(0,0.01)$) -- ($(dqa.north)+(0,-0.01)$);
\draw[gflow] ($(dqa.south)+(0,0.01)$) -- ($(share.north)+(0,-0.01)$);
\draw[rflow] (share.east) -- (multipart.west);
\draw[rflow] (multipart.east) -- ($(proto.west)+(0,0.20)$);
\draw[gflow] (proto.east) -- (enter.west);
\draw[gflow] ($(enter.south)+(0,0.01)$) -- ($(infer.north)+(0,-0.01)$);
\draw[gflow] (infer.east) -- (audit.west);
\draw[gflow] ($(audit.south)+(0,0.01)$) -- ($(reply.north)+(0,-0.01)$);
\draw[bflow] (reply.west) -- ++(-13.55,0) |- (show.east);

% foot notes
\node[footG, minimum width=2.85cm] at (4.10,0.95) {明文止于 Worker 域内};
\node[footR, minimum width=3.20cm] at (7.72,0.95) {跨域仅传 share + 审计摘要};
\node[footO, minimum width=3.05cm] at (11.25,0.95) {SPU 前仍有权威快检};
\node[footB, minimum width=3.00cm] at (14.95,0.95) {最终只回传分类 / 质量 / 审计};

% legend
\draw[tsgreen, line width=1.1pt] (0.35,0.14) -- (1.05,0.14);
\node[anchor=west, font=\fontsize{9.6}{11.5}\selectfont] at (1.18,0.14) {受信任域内明文处理};
\draw[tsred, dashed, line width=1.1pt] (0.35,-0.18) -- (1.05,-0.18);
\node[anchor=west, font=\fontsize{9.6}{11.5}\selectfont] at (1.18,-0.18) {跨边界密态 share / 张量通信};

\end{tikzpicture}
\end{document}
